\begin{table}[H]
\centering
\begin{threeparttable}
\caption{Prueba placebo: tratamiento ficticio en Q2-2021}
\label{tab:tabla5}
\begin{tabular}{lccccc}
\toprule
Lenguaje & ATT Placebo & SE & Significancia & ATT Real (SDID) & SE Real \\
\midrule
C & 0.738 & 0.263 & *** & 1.276 & 0.459 \\
C\# & -0.016 & 0.344 &  & 0.443 & 0.464 \\
C++ & 0.291 & 0.272 &  & 1.131 & 0.293 \\
Go & 0.427 & 0.169 & ** & 0.720 & 0.247 \\
Java & 0.336 & 0.463 &  & 1.725 & 0.835 \\
JavaScript & 3.239 & 1.008 & *** & 8.303 & 2.974 \\
PHP & 0.093 & 0.184 &  & 0.916 & 0.422 \\
Python & 1.694 & 0.639 & *** & 4.525 & 1.057 \\
Ruby & -0.004 & 0.268 &  & 0.808 & 0.282 \\
TypeScript & 1.188 & 0.570 & ** & 3.078 & 0.833 \\
\bottomrule
\end{tabular}
\begin{tablenotes}
\small
\item \textit{Nota.} El tratamiento ficticio se asigna en Q2-2021, período anterior al lanzamiento de ChatGPT. La muestra se restringe a los 11 trimestres pretratamiento (2020-Q1--2022-Q3). Para lenguajes donde el ATT placebo resulta significativo, su magnitud es sustancialmente menor que el ATT real (cociente promedio: 0.38), consistente con la presencia de un efecto de tratamiento genuino por encima de las tendencias pretratamiento.
\item $^{***}p<0.01$, $^{**}p<0.05$, $^{*}p<0.10$.
\end{tablenotes}
\end{threeparttable}
\end{table}
